\section{Dummy-neighborhood corners on fine-cell sides}
\label{sec:area_law_dummy_corners}

Retain the common origin, finite endpoint set and integer width $C_0$,
and fix a lower layer $k_0$. Let $P_{k_0}$ be the indices of the coarse
cells contributing to $N_{k_0}$. Their corners need not be boundary
vertices of the union. The following assertions concern each contributing
cell separately; see \cite[Section~11]{OpenAI2026AreaLaw}.

% Source: 10-geometry.tex, lines 160--191, 200--207 and 299--310.
% Constancy of the opposing dummy region and its label remains separate.

\begin{theorem}[A dummy-cell corner on a whole fine side]
    \label{thm:area_law_dummy_corner_side}
    \lean{TNLean.PEPS.AreaLaw.Geometry.dyadicNeighborhood_corner_on_side}
    \leanok
    \uses{def:area_law_dyadic_layers,def:area_law_cell_corner,
        def:area_law_dyadic_cell_side,def:area_law_fine_layer_indices,
        def:area_law_fine_belt_scales}
    Suppose that $C_0\ge2$, $k\ge k_0$, $z\in F_{k,\ell_k}$ and
    $w\in P_{k_0}$. For $\varepsilon\in\{0,1\}^2$, put
    $x=o+2^{k_0}(w+\varepsilon)$. If $x\in S_{\ell_k,z,s}$, then
    $x=a_{\ell_k,z,s}$ or $x=b_{\ell_k,z,s}$.
\end{theorem}
\begin{proof}\leanok
    \uses{thm:area_law_dummy_contact_index,thm:area_law_fine_cell_containment,
        thm:area_law_fine_belt_scales,thm:area_law_cell_fan_mark_vertices,
        thm:area_law_belt_marks_closed_cell,thm:area_law_dyadic_closures,
        thm:area_law_dyadic_cell_distance,
        def:area_law_affine_mesh,thm:area_law_affine_mesh_point_separation}
    The corner belongs to $\overline{N_{k_0}}$. The whole side lies in
    the convex square $\overline{Q_{k,z}}\subseteq\overline{D_k}$,
    so dummy-contact localization gives $k=k_0$. Put $q=2^{\ell_{k_0}}$.
    Since $\ell_{k_0}\le k_0$, the corner and both side endpoints belong
    to $o+q\mathbb Z^2$. If $x$ differed from both endpoints $a,b$,
    mesh separation would give
    \begin{align}
        2q\le\|a-x\|_\infty+\|x-b\|_\infty
        =\|a-b\|_\infty\le q.
    \end{align}
    The equality follows from $x\in[a,b]$, and the last inequality
    is the cell's sup-diameter bound. This contradicts $q>0$.
\end{proof}

\begin{theorem}[A dummy-cell corner on an elementary fine side]
    \label{thm:area_law_dummy_corner_elementary_side}
    \lean{TNLean.PEPS.AreaLaw.Geometry.dyadicNeighborhood_corner_on_elementarySide}
    \leanok
    \uses{def:area_law_dyadic_layers,def:area_law_cell_corner,
        def:area_law_cell_fan,def:area_law_fine_layer_indices,
        def:area_law_fine_belt_scales}
    Suppose that $C_0\ge2$, $k\ge k_0$, $z\in F_{k,\ell_k}$ and
    $w\in P_{k_0}$. Choose any optional midpoint subdivision of
    $\overline{Q_{k,z}}$ and one of its elementary segments $[u,v]$.
    For $\varepsilon\in\{0,1\}^2$, if the corner
    $x=o+2^{k_0}(w+\varepsilon)$ of the contributing dummy cell belongs
    to $[u,v]$, then $x=u$ or $x=v$.
\end{theorem}
\begin{proof}\leanok
    \uses{thm:area_law_dummy_corner_side,thm:area_law_elementary_side_geometry}
    The elementary segment belongs to its whole side $[a,b]$, so
    $x=a$ or $x=b$. For an undivided side this is the desired conclusion.
    On $[a,(a+b)/2]$, the endpoint $b$ cannot occur unless $a=b$;
    on $[(a+b)/2,b]$, the endpoint $a$ cannot occur unless $a=b$.
    If the endpoints coincide, the conclusion holds directly. Thus
    every possible case gives one of the elementary endpoints.
\end{proof}
